Gravity and magnetic anomaly maps offer a position reference that cannot be jammed, and a vehicle-mounted smartphone already carries both an accelerometer and a magnetometer. This paper asks whether those readings, or the readings of comparably priced dedicated sensors, reduce the error of a smartphone-grade inertial navigation system (INS) when GNSS fixes arrive only once per minute with correlated errors. An extended Kalman filter (EKF), an unscented Kalman filter (UKF), and a Rao-Blackwellized particle filter (RBPF) after Canciani and Raquet are evaluated on 27 drives totalling \numdatasetkm~km and \numdatasethours~h, each aided run paired with the same filter's unaided run on an identical, seeded measurement stream. The paper corrects and supersedes two earlier conference papers by the authors, whose reported gains arose from unit errors in the anomaly measurement and from comparisons against diverged baselines. With the corrected models, the smartphone's own readings provide no measurable benefit to either Kalman filter (median RMSE ratio 1.000 to 1.001), because the operating system pins the gravity norm to a device constant and the magnetometer is dominated by the vehicle's field, giving signal-to-noise ratios of 0.13 and 0.01 against the maps. Replacing only those readings with a simulated RM3100 magnetometer reduces the median error by about 10\% (ratios \numdedicatedekfmagratio{} for the EKF and \numdedicatedukfmagratio{} for the UKF, improving \numdedicatedekfmagbetter{} and \numdedicatedukfmagbetter{} of 27 drives), whereas a simulated ADXL355 used as a gravimeter does not help. The RBPF loses track under the 60-s GNSS schedule with or without aiding. The quality of the anomaly measurement relative to the map signal, rather than the estimator or the inertial sensor, decides whether aiding helps.
